\documentclass[12pt]{article}
\usepackage[T1]{fontenc}
\usepackage[utf8]{inputenc}
\usepackage{lmodern}
\usepackage{textcomp}
\usepackage{enumitem}
\usepackage{geometry}
\geometry{margin=1in}
\begin{document}
\begin{center}
Answer Key - Business Administration - Undergraduate Level Assessment 5 \
\small (Answers are ordered exactly as in the question paper.)
\end{center}

\section*{Multiple Choice}
\begin{enumerate}[label=\arabic*.]
\item \textbf{Answer:} D
\\ \textit{Options:} A) Programmed decision; B) Routine decision; C) Management decision; D) Non-programmed decision
\\ \textit{Note:} An investment decision is considered a non-programmed decision because it involves unique, complex situations that require judgment and discretion rather than relying on established rules or procedures.
\item \textbf{Answer:} B
\\ \textit{Options:} A) Individual branding.; B) Family branding.; C) Corporate brands.; D) Manufacturer brand.
\\ \textit{Note:} Family branding is the marketing strategy where a company uses its brand name across multiple products, thereby leveraging the established reputation of the brand to promote all items under that name, as seen with companies like Microsoft, Heinz, and Kellogg's.
\item \textbf{Answer:} D
\\ \textit{Options:} A) Peer recognition; B) Promotion; C) Greater freedom; D) Economic reward
\\ \textit{Note:} Workers are primarily motivated by economic reward according to the instrumental approach because it emphasizes the role of financial incentives as a key driver of employee behavior and performance in the workplace.
\item \textbf{Answer:} B
\\ \textit{Options:} A) 3,4; B) 1,3; C) 2,3; D) 4,1
\\ \textit{Note:} The correct answer is based on the understanding that certain business practices or regulations can be perceived as irrational and lacking a formal codification, which can lead to inconsistencies and challenges in their application within the field of business administration.
\item \textbf{Answer:} D
\\ \textit{Options:} A) Cognition.; B) Perception.; C) Learning.; D) Personality.
\\ \textit{Note:} Personality is not a fundamental process in consumer learning and decision-making; instead, it is a trait that influences behavior but does not directly explain the cognitive processes involved in understanding products and brands.
\end{enumerate}

\section*{True / False}
\begin{enumerate}[label=\arabic*.]
\setcounter{enumi}{5}
\item \textbf{Answer:} False
\\ \textit{Options:} A) True; B) False
\\ \textit{Note:} Luxury products are often high-involvement purchases that require consumers to conduct thorough research and update their knowledge about brand reputation, quality, and market trends before making a purchase decision.
\item \textbf{Answer:} False
\\ \textit{Options:} A) True; B) False
\\ \textit{Note:} Encoding refers to the process of converting thoughts or ideas into symbols or messages by the sender, not the responses provided by the receivers.
\item \textbf{Answer:} True
\\ \textit{Options:} A) True; B) False
\\ \textit{Note:} A hierarchical structure often involves multiple layers of management and formal procedures, which can slow down decision-making and hinder the organization's ability to adapt swiftly to changes in the business environment.
\item \textbf{Answer:} False
\\ \textit{Options:} A) True; B) False
\\ \textit{Note:} Task culture is designed to foster flexibility and adaptability in response to changing project needs, which can sometimes undermine stability and efficiency compared to more hierarchical or bureaucratic cultures.
\item \textbf{Answer:} True
\\ \textit{Options:} A) True; B) False
\\ \textit{Note:} The statement is true because diffusion refers to the process by which an innovation is communicated and adopted within a market or social system over time.
\end{enumerate}

\section*{Long Form Response}
\begin{enumerate}[label=\arabic*.]
\setcounter{enumi}{10}
\item \textbf{Answer:} Issues management and crisis management both address acute risks, but they differ in their approaches and timing. Issues management is proactive, focusing on identifying and addressing potential problems before they escalate into crises. For example, a company may implement a stakeholder engagement strategy to mitigate negative perceptions related to environmental practices. In contrast, crisis management is reactive, dealing with situations that have already escalated into crises, such as a product recall due to safety concerns. Both disciplines aim to protect the organization's reputation and ensure business continuity, but they operate at different stages of risk management, highlighting their distinct yet complementary roles.
\\ \textit{Note:} The answer correctly identifies that issues management is proactive in preventing potential problems, while crisis management is reactive in addressing situations that have already escalated, using relevant examples to illustrate these distinctions.
\item \textbf{Answer:} Fayol's 14 principles of management emphasize the importance of structured organization, teamwork, and collective efforts in achieving business goals. These principles advocate for unity of command, division of work, and discipline, which collectively foster a collaborative environment. Individualism, which prioritizes personal interests and autonomy, conflicts with these principles as it can lead to fragmented efforts and reduced coherence within a team. Effective management practices require alignment with Fayol's principles to ensure that all employees work towards a common objective, thus enhancing organizational efficiency and cohesion. Therefore, while individual contributions are valuable, they must be harmonized with a collective approach to management to achieve optimal results.
\\ \textit{Note:} correct because it articulates how Fayol's principles promote teamwork and organizational coherence, contrasting with individualism's focus on personal interests, which can undermine collaborative efforts essential for effective management.
\end{enumerate}

\vfill
\noindent\textit{End of Answer Key}
\end{document}